\section*{Chapter \ref{ch:b3:environmental-toxicology} --- Environmental Chemistry and Toxicology}

\begin{solution}{exo:b3:environmental-toxicology:1}
A homonuclear diatomic has a single vibration that keeps the dipole moment zero,
and argon has no vibration: none can absorb infrared radiation. $\ce{CO2}$ absorbs through
its bending mode (and its antisymmetric stretch), which create an oscillating
dipole; the bend lies in the middle of the Earth's emission.
\end{solution}

\begin{solution}{exo:b3:environmental-toxicology:2}
$10^{0.1} = 1.26$: $[\ce{H+}]$ rises by 26\,\%.
\end{solution}

\begin{solution}{exo:b3:environmental-toxicology:3}
$(2.0 + 0.50)\ \unit{mmol/L} \times \qty{100.1}{mg/mmol} = \qty{250}{mg/L}$ of $\ce{CaCO3}$.
\end{solution}

\begin{solution}{exo:b3:environmental-toxicology:4}
N: $0.60/14.0 = \qty{0.043}{mmol/L}$; P: $0.010/31.0 = \qty{3.2e-4}{mmol/L}$; N\,:\,P $= 133$, far above 16: phosphorus runs out
first and limits growth.
\end{solution}

\begin{solution}{exo:b3:environmental-toxicology:5}
$L_0 = \mathrm{BOD_5}/(1 - \eu^{-0.23 \times 5}) = 170/0.683 = \qty{250}{mg/L}$.
\end{solution}

\begin{solution}{exo:b3:environmental-toxicology:6}
$\qty{4.40e-3}{L} \times \qty{0.0200}{mol/L} = \qty{8.80e-5}{mol}$ in \qty{0.1000}{L}: \omterm{def:b3:environmental-toxicology:alkalinity}{alkalinity} $\qty{8.80e-4}{mol/L}$; as $\ce{HCO3-}$ (\qty{61.0}{g/mol}),
\qty{53.7}{mg/L}.
\end{solution}

\begin{solution}{exo:b3:environmental-toxicology:7}
$K_d = 0.020 \times 300 = \qty{6.0}{L/kg}$. Sorbed/total $= K_dm_s/(K_dm_s + V_w) = 9.0/(9.0 + 0.30) = 0.97$: 97\,\% sorbed.
\end{solution}

\begin{solution}{exo:b3:environmental-toxicology:8}
Benzene: $\log\mathrm{BCF} = 0.85 \times 2.13 - 0.70 = 1.11$, BCF about 13. PCB: $0.85 \times 6.67 - 0.70 = 4.97$, BCF about
\num{9e4}: it concentrates in fish some seven thousand times more than benzene; at such
$K_{ow}$ the regression is near its limit of validity.
\end{solution}

\begin{solution}{exo:b3:environmental-toxicology:9}
$\qty{192}{mg/kg} \times \qty{70}{kg} = \qty{13}{g}$, about 150 cups. Species differ in absorption and metabolism,
and an $\mathrm{LD_{50}}$ is not a safe threshold; harmful effects begin far below it.
\end{solution}

\begin{solution}{exo:b3:environmental-toxicology:10}
PNEC $= \qty{0.50}{mg/L}/1000 = \qty{0.50}{\micro g/L}$; $\mathrm{RQ} = 0.20/0.50 = 0.40$, below 1: no concern at this exposure.
\end{solution}

\begin{solution}{exo:b3:environmental-toxicology:11}
Plankton/water: $\qty{0.015}{mg/kg}/\qty{2.0e-6}{mg/L} = \qty{7.5e3}{L/kg}$ (bioconcentration). Fish/plankton: $0.30/0.015 = 20$;
bird/fish: $6.0/0.30 = 20$ (\omterm{def:b3:environmental-toxicology:bio}{biomagnification} at each step); bird/water: $\num{3e6}$.
\end{solution}

\begin{solution}{exo:b3:environmental-toxicology:12}
$\qty{1.0e9}{g}/\qty{64.1}{g/mol} = \qty{1.56e7}{mol}$ of $\ce{SO2}$, giving as many moles of $\ce{H2SO4}$: $\qty{1.53e9}{g}$, about \qty{1500}{t}. It
releases $\qty{3.12e7}{mol}$ of $\ce{H+}$; at pH 4.0, $[\ce{H+}] = \qty{1.0e-4}{mol/L}$: $\qty{3.1e11}{L}$ of rain.
\end{solution}

\begin{solution}{pb:b3:environmental-toxicology:1}
\textbf{1.} It is ten thousand times more soluble in octanol than in water: hydrophobic,
it will sorb on organic matter and accumulate in fat.

\textbf{2.} $K_{oc} = 0.41 \times 10^4 = \qty{4100}{L/kg}$; $K_{sw} = f_{oc}K_{oc} = 0.040 \times 4100 = \qty{164}{L/kg}$.

\textbf{3.} $K_{fw} = 0.048 \times 10^4 = \qty{480}{L/kg}$.

\textbf{4.} No: at equilibrium the air holds $10^{-4}$ of the water concentration.

\textbf{5.} The neutral, hydrophobic molecule dissolves in the organic matter, as in
octanol; mineral surfaces are polar and covered with water.

\textbf{6.} The sediment, whose capacity, $K_{sw}m_s$, is largest.

\textbf{7.} $V_w = \qty{1.0e9}{L}$; $K_{aw}V_a = \qty{1.0e8}{L}$; $K_{sw}m_s = \qty{1.64e10}{L}$; $K_{fw}m_f = \qty{4.8e6}{L}$.

\textbf{8.} $C_w = \qty{1.0e8}{mg}/\qty{1.75e10}{L} = \qty{5.7e-3}{mg/L}$, \qty{5.7}{\micro g/L}.

\textbf{9.} Water \qty{5.7}{kg}, air \qty{0.57}{kg}, sediment \qty{94}{kg}, fish \qty{0.027}{kg}.

\textbf{10.} Air $\qty{5.7e-7}{mg/L}$; sediment \qty{0.94}{mg/kg}; fish \qty{2.7}{mg/kg}.

\textbf{11.} Equilibrium between all compartments at once, with no degradation, no
inflow or outflow, and well-mixed compartments.

\textbf{12.} With degradation and flows at steady state (level II) the inventory is set by
the rates of input and loss; with transfer rates between compartments (level III)
the concentrations are no longer in equilibrium ratios, and the compartment that
receives the input is enriched.

\textbf{13.} $k = \ln 2/\qty{60}{d} = \qty{0.0116}{d^{-1}}$.

\textbf{14.} $\eu^{-0.0116 \times 365} = 0.0145$: 1.5\,\% left.

\textbf{15.} $\qty{5.7}{\micro g/L} \times 0.0145 = \qty{0.083}{\micro g/L}$.

\textbf{16.} Degradation is often slower in cold, anoxic sediment, and the sorbed
pesticide is released back to the water slowly, keeping it present after the
water has been cleared.

\textbf{17.} Bonds that resist hydrolysis, oxidation and microbial attack (aromatic
C--Cl, for example), low water solubility, and sorption that protects it from
degraders.

\textbf{18.} PNEC $= \qty{50}{\micro g/L}/10 = \qty{5.0}{\micro g/L}$.

\textbf{19.} $\mathrm{RQ} = 5.7/5.0 = 1.1$.

\textbf{20.} Just above 1: a risk is indicated, and the assessment must be refined (better
toxicity data, measured concentrations) or the exposure reduced.

\textbf{21.} \qty{2.7}{mg/kg}, almost three times the \qty{1.0}{mg/kg} threshold: the fish should not be eaten
until the concentration falls.

\textbf{22.} Concentrations in water, sediment and fish over time and at several points,
to check the model and follow the decline.

\textbf{23.} Less pesticide applied, application away from rain and from the shore,
buffer strips that hold runoff, or a less persistent alternative.

\textbf{24.} It divides the measured no-effect level to cover the gap between a few
laboratory species and a whole ecosystem; a factor of 10 instead of 1000 changes the
PNEC, and the verdict, a hundredfold.

\textbf{25.} The \omterm{def:b3:environmental-toxicology:ecotoxicology}{risk quotient} PEC/PNEC is about 1.1, just above 1.
\end{solution}
